\scanpage{36}
(evident from (8)), which \deleted{\unsure{is defined}}\inserted{can be described} as the set of $(p,\sigma)$ $(p\in\hat{\mathbb Z}^*,\ \sigma\in\unsure{\mathfrak S_3|\mathcal S_3})$ such that one has
\[
(13)\qquad l_{\sigma(0)}^p\ l_{\sigma(1)}^p\ l_{\sigma(\infty)}^p=1
\]
When one restricts oneself to $S\mathcal E$ instead of $S\widehat{\widehat{\mathcal E}}$, so that one must have $p\in\mathbb Z^*=\{\pm1\}$, \deleted{one finds that the \unsure{2}} i.e. the set of $(p,\sigma)$ (\unsure{abstracting from|forgetting} the condition (13)) is none other than $\deleted{\unsure{\mathbb Z}}\inserted{\{\pm1\}}\times\unsure{\mathfrak S_3|\mathcal S_3}\simeq\unsure{\mathfrak S_3|\mathcal S_3}\times\mathbb Z_2\simeq\Gamma$, one sees that $(p,\sigma)$ satisfies (13) if and only if $(\sigma,p)\in\Gamma_P$. Thus, $\tilde\Gamma_P\subset S\mathcal E\subset S\widehat{\widehat{\mathcal E}}$ is characterized as the subgroup \inserted{of $S\mathcal E$} of automorphisms of $\pi$ \deleted{\unsure{or|in} $\widehat{\widehat\pi}$} which leave $L_0,L_1,L_\infty$ invariant, equipped with the \deleted{multipliers}\inserted{``rectifiers''} $g_0\unsure{=|\hat{=}|\overset{\cdot}{=}}1,g_1=1,g_\infty=1$.

Probably it is true that the same result is valid in $S\widehat{\widehat{\mathcal E}}$, i.e. that the relation (13) necessarily implies $p\in\{\pm1\}$. It amounts to the same thing to say \inserted{(up \unsure{:|to} correcting by a $\tilde\sigma_i$ --- \unsure{or} $\sigma\in\Gamma_P$)} that the relation
\[
(14)\qquad :\ l_\infty^p\ l_1^p\ l_0^p=1
\]
(for a fixed $p\in\hat{\mathbb Z}^*$) implies $p=1$.

Exercise!
\mspage{35}
